% =============================================================================
% Draft supplementary table for Q3 -- PINN/PINN-k across environmental feature
% sets. Companion to Table~\ref{tab:results-q3} (which stays fixed at Set3 for
% cross-model comparability -- see chat discussion 2026-08-23). This table
% carries the "does more environment help" finding on its own, without
% touching the main model-comparison table.
% Numbers from temp_results_pinn/RESULTS_TABLE.md, section 5.
% =============================================================================

\begin{table}[!htbp]
\centering
\small
\setlength{\tabcolsep}{6pt}
\renewcommand{\arraystretch}{1.2}
\begin{tabular}{lrr}
\toprule
Feature set & PINN $R^2{\uparrow}$ & PINN-$k$ $R^2{\uparrow}$ \\
\midrule
No environment ($\text{pinn\_noenv}$) & $0.573$ & $0.573$ \\
Set2 (top-10 features) & $0.627{\pm}0.023$ & $0.622{\pm}0.017$ \\
Set3 (VIF-gated, headline -- Table~\ref{tab:results-q3}) & $\mathbf{0.631{\pm}0.027}$ & $0.618{\pm}0.023$ \\
Set4 (broadest, all VIF-gated) & $0.618{\pm}0.029$ & $\mathbf{0.620{\pm}0.020}$ \\
\bottomrule
\end{tabular}
\caption{\small{PINN and PINN-$k$ accuracy (Set3, spatial\_block\_kfold, per-fold mean $\pm$ SD
across five folds, 4-survey) across environmental feature-set breadth, from no environment
information at all up to the broadest VIF-screened set. Bold marks each column's own best
mean, shown for transparency, not as a claim that the difference is reliable --
Set2/Set3/Set4 sit within one fold-to-fold SD of each other for both variants (a
difference smaller than the $\sim$0.02--0.03 noise floor), so no set is a demonstrated
winner among the three for either variant. The one large, reliable difference is between
having environment information at all and not having it (No environment $\to$ Set2:
$+0.054$ PINN, $+0.049$ PINN-$k$, several times the noise floor).}}
\label{tab:results-q3-env-sweep}
\end{table}

% -----------------------------------------------------------------------------
% Draft accompanying prose
% -----------------------------------------------------------------------------

\paragraph*{Does more environmental information help PINN?}
\textbf{Environmental conditioning delivers a real, substantial gain once -- moving from no
environment to any environment -- but does not improve further with a broader feature list,
and Set3's selection as the headline set (Table~\ref{tab:results-q3}) is not itself a
demonstrated best-performer for PINN-$k$.}
Table~\ref{tab:results-q3-env-sweep} compares PINN and PINN-$k$ with no environmental
information at all against three environmental feature sets of increasing breadth. The jump
from no environment to any environment is large and reliable ($+0.054$ $R^2$ for PINN,
$+0.049$ for PINN-$k$, several times the $\sim$0.02--0.03 fold-to-fold noise floor). Beyond
that, feature-set breadth does not matter: Set2, Set3, and Set4 sit within one SD of each
other for both variants, so no set is a statistically demonstrated winner among the three.
Set3 is used as the headline feature set throughout this chapter for consistency with every
other model in Table~\ref{tab:results-q3} (Linear, RF, XGBoost, DNN all also report at
Set3) -- a defensible choice for cross-model comparability, but not, on this evidence, a
claim that Set3 specifically maximises PINN-$k$'s accuracy.
